\documentclass[11pt,a4paper]{beamer}
\usetheme{Madrid}
\usepackage[utf8]{inputenc}
\usepackage{amsmath}
\usepackage{amsfonts}
\usepackage{amssymb}
\usepackage{graphicx}
\usepackage{tikz}
\usepackage{booktabs}
\usepackage{array}
\usepackage{colortbl}
\usepackage{multicol}

\usefonttheme{professionalfonts}

\setbeamerfont{normal text}{size=\tiny}
\setbeamerfont{itemize/enumerate body}{size=\tiny}
\setbeamerfont{itemize/enumerate subbody}{size=\tiny}
\setbeamerfont{block body}{size=\tiny}
\setbeamerfont{block title}{size=\scriptsize}
\setbeamerfont{frametitle}{size=\small}
\setbeamerfont{title}{size=\small}
\setbeamerfont{subtitle}{size=\footnotesize}
\setbeamerfont{author}{size=\tiny}
\setbeamerfont{date}{size=\tiny}

\setlength{\leftmargini}{0.3em}
\setlength{\leftmarginii}{0.3em}
\setlength{\labelsep}{0.1em}
\setlength{\itemsep}{0.02em}
\setlength{\parskip}{0.02em}

\title[CAMS Exam Preparation]{CAMS Exam Preparation}
\subtitle{Broker-Dealer \& Investment Bank AML Risks}
\author{Prime Brokerage | Insider Trading | Cross-Border | IPOs}
\date{}

\setbeamercolor{title}{fg=white,bg=blue!70!black}
\setbeamercolor{frametitle}{fg=white,bg=blue!70!black}
\setbeamertemplate{navigation symbols}{}

\begin{document}

% TITLE SLIDE
\begin{frame}
\titlepage
\centering
\tiny 15 Topics | Broker-Dealer AML Risks | Natural Explanations
\end{frame}

% ============================================================
% SLIDE 1 - BROKER-DEALER AML OBLIGATIONS
% ============================================================
\begin{frame}{1. Broker-Dealer AML Obligations}
\tiny
\noindent\textbf{Introduction to the topic:} Broker-dealers execute securities trades for customers. Unlike banks, they may not hold customer deposits for long periods. However, they are fully subject to AML regulations under the Bank Secrecy Act and FINRA Rule 3310. They must establish AML programs, designate a compliance officer, conduct customer due diligence (including identifying beneficial owners), monitor transactions for suspicious activity, and file SARs. The fact that a broker-dealer does not hold customer funds does not reduce its AML obligations — trade execution itself creates money laundering risk through layering and integration.

\vspace{0.1cm}
\noindent\textbf{Type of violation or warning signs:} A broker-dealer that processes high trade volume but has no AML program or has never filed a SAR is a warning sign. So is a broker-dealer that relies entirely on an introducing broker for customer identification without independent verification. Other red flags include accepting instructions from unverified third parties, executing trades for shell companies without UBO identification, and processing rapid in-and-out trading that generates losses (potential layering).

\vspace{0.1cm}
\noindent\textbf{Correct answer approach:} For exam questions about broker-dealer obligations, the correct answer will state that broker-dealers must have a full AML program including CDD, transaction monitoring, and SAR filing. Eliminate any answer suggesting broker-dealers are exempt or have reduced obligations. Remember that trade execution alone creates money laundering risk, and every registered broker-dealer is subject to the same AML rules as banks.
\end{frame}

% ============================================================
% SLIDE 2 - OMNIBUS ACCOUNTS & INTRODUCING BROKERS
% ============================================================
\begin{frame}{2. Omnibus Accounts and Introducing Brokers}
\tiny
\noindent\textbf{Introduction to the topic:} An omnibus account is a single account held by a carrying broker that contains the aggregated positions of multiple customers of an introducing broker. The carrying broker sees only the introducing broker's name, not the underlying customers. This structure creates significant AML risk because the carrying broker cannot identify who is actually trading. FINRA and SEC rules require carrying brokers to have agreements with introducing brokers that ensure customer identification information is provided upon request. Without this, the carrying broker is effectively blind.

\vspace{0.1cm}
\noindent\textbf{Type of violation or warning signs:} The carrying broker has no customer identification information for omnibus account transactions. The introducing broker refuses to provide underlying customer names citing “proprietary information.” The carrying broker processes trades for the omnibus account without any CDD on the ultimate traders. Regulators have fined carrying brokers for failing to obtain customer information from introducing brokers. Another warning sign is high-volume trading through an omnibus account from a high-risk jurisdiction.

\vspace{0.1cm}
\noindent\textbf{Correct answer approach:} The correct answer will state that the carrying broker must have agreements in place to obtain customer identification information from introducing brokers. If the introducing broker refuses to provide underlying customer names, the carrying broker must terminate the relationship. Eliminate answers that suggest the carrying broker has no obligation or that omnibus accounts are automatically acceptable. Remember: the carrying broker is responsible for knowing its customers, even when they trade through an introducing broker.
\end{frame}

% ============================================================
% SLIDE 3 - INSIDER TRADING AS MONEY LAUNDERING
% ============================================================
\begin{frame}{3. Insider Trading as a Predicate Offense}
\tiny
\noindent\textbf{Introduction to the topic:} Insider trading is not just a securities violation — it is a predicate offense for money laundering in most jurisdictions. When a person profits from insider trading, those profits are “criminal proceeds” under money laundering statutes. Using a brokerage account to trade on inside information, then moving the proceeds to another account or jurisdiction, constitutes money laundering. The money laundering charge can carry longer sentences than the insider trading violation itself.

\vspace{0.1cm}
\noindent\textbf{Type of violation or warning signs:} A customer with no history of profitable trading suddenly has an unusually high win rate on specific stocks. Trades occur just before major corporate announcements (earnings, mergers, FDA approvals). The customer cannot explain the source of their trading information. The customer moves trading profits immediately to offshore accounts or converts to cryptocurrency. Multiple accounts trade the same stock just before the same announcement from different introducing brokers but with similar patterns.

\vspace{0.1cm}
\noindent\textbf{Correct answer approach:} The correct answer will identify insider trading profits as criminal proceeds subject to money laundering laws. For exam questions, look for the answer that states insider trading is a predicate offense and that brokerage firms must file SARs for suspected insider trading. Eliminate answers that treat insider trading as only a securities violation or suggest that AML obligations do not apply. Remember: any SAR filed for insider trading should also consider money laundering charges.
\end{frame}

% ============================================================
% SLIDE 4 - MARKET MANIPULATION (PUMP AND DUMP)
% ============================================================
\begin{frame}{4. Market Manipulation (Pump and Dump)}
\tiny
\noindent\textbf{Introduction to the topic:} Pump and dump schemes involve artificially inflating the price of a stock (often a low-float penny stock) through false or misleading statements, then selling shares at the inflated price. The proceeds are criminal proceeds subject to money laundering laws. Brokers that execute these trades — especially those that ignore red flags about the issuer or the trading pattern — can be liable for facilitating money laundering. The manipulation often involves multiple accounts controlled by the same person.

\vspace{0.1cm}
\noindent\textbf{Type of violation or warning signs:} Unusual volume and price movement in a low-priced stock with no news. The same customer account sells large blocks after the price spikes. Multiple accounts from different introducing brokers sell the same stock at the same time. The issuer has no legitimate business operations (shell company). The customer cannot explain the source of the trading idea or the basis for the stock's valuation. Rapid movement of sale proceeds to offshore accounts or cryptocurrency.

\vspace{0.1cm}
\noindent\textbf{Correct answer approach:} The correct answer will identify that pump-and-dump proceeds are laundered through brokerage accounts and that brokers must monitor for manipulation patterns. For exam questions, look for answers that require SAR filing for suspected manipulation and enhanced due diligence on penny stock trading. Eliminate answers that suggest brokers have no responsibility for monitoring trading patterns or that manipulation is only an SEC issue, not an AML issue.
\end{frame}

% ============================================================
% SLIDE 5 - CROSS-BORDER BROKERAGE FLOWS
% ============================================================
\begin{frame}{5. Cross-Border Brokerage Flows}
\tiny
\noindent\textbf{Introduction to the topic:} Cross-border brokerage flows involve moving funds between brokerage accounts in different countries. This is a common money laundering technique for layering. Funds may move from a brokerage account in Country A to an account in Country B, then to Country C, then back to Country A, with no economic purpose. Each transfer adds another layer of obfuscation. Brokers must monitor for these patterns, especially when involving high-risk jurisdictions or jurisdictions with bank secrecy laws.

\vspace{0.1cm}
\noindent\textbf{Type of violation or warning signs:} Funds move from a brokerage account to another brokerage account in a different country, then to a third country, with no investment activity at each stop. The customer cannot explain the business purpose of the transfers. The transfers involve jurisdictions known for weak AML enforcement or bank secrecy. The customer requests that trades be executed through a broker in a high-risk jurisdiction rather than a local broker. Rapid movement of funds just below reporting thresholds.

\vspace{0.1cm}
\noindent\textbf{Correct answer approach:} The correct answer will state that cross-border movements without economic purpose indicate layering and require SAR filing. For exam questions, look for answers that require the broker to ask for the business purpose of each transfer and file SAR if the explanation is inadequate. Eliminate answers that suggest cross-border transfers are normal and require no scrutiny. Remember: the number of hops matters — three or more jurisdictions with no investment activity is almost certainly layering.
\end{frame}

% ============================================================
% SLIDE 6 - PRIME BROKERAGE RISKS
% ============================================================
\begin{frame}{6. Prime Brokerage Risks}
\tiny
\noindent\textbf{Introduction to the topic:} Prime brokers provide clearing, custody, and financing services to hedge funds and other professional trading firms. The prime broker often does not know the underlying investors in the hedge fund — it knows only the hedge fund itself. This creates AML risk because the prime broker may be processing trades for a fund whose investors include PEPs, sanctioned persons, or criminals. The prime broker must perform due diligence on the hedge fund's investor base, not just on the fund itself.

\vspace{0.1cm}
\noindent\textbf{Type of violation or warning signs:} The hedge fund refuses to provide information about its investors. The fund is domiciled in a secrecy jurisdiction (Cayman, BVI, Bermuda). The fund has no independent administrator or auditor. The fund's investors include foreign PEPs or persons from high-risk jurisdictions. The fund engages in highly complex, rapid trading that generates consistent small losses (potential layering). The fund uses multiple prime brokers to fragment its trading.

\vspace{0.1cm}
\noindent\textbf{Correct answer approach:} The correct answer will state that prime brokers must perform due diligence on the hedge fund's underlying investors to identify PEPs, sanctions exposure, and source of funds. If the hedge fund refuses to provide investor information, the prime broker must decline or terminate the relationship. Eliminate answers that suggest the prime broker's only obligation is to the hedge fund entity itself. Remember: the prime broker must know who ultimately controls the funds being traded.
\end{frame}

% ============================================================
% SLIDE 7 - SPONSORED ACCESS (NAKED ACCESS)
% ============================================================
\begin{frame}{7. Sponsored Access (Naked Access)}
\tiny
\noindent\textbf{Introduction to the topic:} Sponsored access (also called naked access) allows a customer to trade directly on an exchange using the broker's market participant identifier, without pre-trade risk controls or AML screening by the broker. The customer's trades appear to the exchange as if they come from the broker. This practice is highly risky because the broker has no real-time ability to monitor or block suspicious trading. The SEC has restricted naked access, requiring brokers to have risk controls in place.

\vspace{0.1cm}
\noindent\textbf{Type of violation or warning signs:} The broker provides sponsored access to a customer without pre-trade risk controls. The customer is located in a high-risk jurisdiction. The customer's trading volume is extremely high relative to their stated net worth. The broker cannot monitor customer trades in real time. The customer refuses to provide information about its underlying traders. The customer engages in market manipulation or insider trading using the broker's identifier, making it appear that the broker itself is trading.

\vspace{0.1cm}
\noindent\textbf{Correct answer approach:} The correct answer will state that sponsored access without pre-trade risk controls and AML screening is prohibited or requires significant mitigation. For exam questions, look for answers that require the broker to implement real-time monitoring and risk controls for any sponsored access customer. Eliminate answers that suggest naked access is acceptable for certain customers. Remember: if the customer trades under the broker's identifier, the broker is legally responsible for those trades.
\end{frame}

% ============================================================
% SLIDE 8 - IPOs AND PIPE TRANSACTIONS
% ============================================================
\begin{frame}{8. IPOs and PIPE Transactions}
\tiny
\noindent\textbf{Introduction to the topic:} Initial Public Offerings (IPOs) and Private Investments in Public Equity (PIPE) transactions involve the sale of securities to investors before they trade publicly. Criminals use IPOs and PIPEs to launder money by purchasing shares with illicit funds, then selling them after the lock-up period for clean proceeds. The difficulty is verifying the source of funds for these investments — investors may use shell companies, offshore trusts, or nominee accounts to hide their identity.

\vspace{0.1cm}
\noindent\textbf{Type of violation or warning signs:} An IPO investor uses a shell company or trust in a secrecy jurisdiction to purchase shares. The investor cannot provide documentation for the source of funds. The investor is a foreign PEP or has close associates who are PEPs. The investor purchases a large allocation in an IPO but has no history of public equity investing. The investor sells the shares immediately after the lock-up period ends, regardless of price. The investor uses multiple nominee accounts to purchase the same IPO.

\vspace{0.1cm}
\noindent\textbf{Correct answer approach:} The correct answer will state that brokers and underwriters must verify the source of funds for IPO and PIPE investors, especially those using complex structures or from high-risk jurisdictions. If source of funds cannot be verified, the investor should be rejected. Eliminate answers that suggest institutional investors are automatically acceptable or that source of funds verification is optional. Remember: the lock-up period does not cleanse the funds — the original source remains relevant.
\end{frame}

% ============================================================
% SLIDE 9 - SHELL COMPANY LISTINGS
% ============================================================
\begin{frame}{9. Shell Company Listings (Reverse Mergers)}
\tiny
\noindent\textbf{Introduction to the topic:} Shell companies (public companies with no operations) are sometimes used for reverse mergers, allowing private companies to become public without a traditional IPO. Criminals use shell companies to launder money by acquiring control of the shell, injecting illicit funds into the shell through fake transactions, and then selling shares to the public. The shell company may have no legitimate business, inadequate controls, and a management team with criminal backgrounds.

\vspace{0.1cm}
\noindent\textbf{Type of violation or warning signs:} A shell company with no operations suddenly reports huge revenues from an unexplained source. The shell company's management has backgrounds in penny stock promotion or have been involved in prior SEC enforcement actions. The shell company is domiciled in a secrecy jurisdiction. The company issues shares to entities in high-risk jurisdictions. Trading volume in the stock is unusually high with no corresponding news. The company refuses to provide audited financial statements.

\vspace{0.1cm}
\noindent\textbf{Correct answer approach:} The correct answer will state that brokers must conduct enhanced due diligence on shell company listings, including verifying the source of the company's assets and the background of management. If red flags exist, the broker should restrict or prohibit trading in the stock. Eliminate answers that suggest shell companies are legitimate simply because they are public. Remember: many shell companies are vehicles for money laundering and fraud.
\end{frame}

% ============================================================
% SLIDE 10 - BROKER-CUSTODIAN CONFLICTS
% ============================================================
\begin{frame}{10. Broker-Custodian Conflicts of Interest}
\tiny
\noindent\textbf{Introduction to the topic:} Many brokerage firms also act as custodians for customer assets. This creates a potential conflict of interest: the firm may be reluctant to file SARs or exit high-risk customers because those customers generate custody fees. Regulators examine whether custody revenue influences AML decisions. A firm that retains a high-risk customer primarily for custody fees is violating AML obligations. The firm must have independent AML decision-making separate from revenue considerations.

\vspace{0.1cm}
\noindent\textbf{Type of violation or warning signs:} The firm's largest custody customers are also its highest-risk customers (PEPs, offshore structures). The firm has never filed a SAR on these customers despite red flags. Compliance recommended terminating a customer but senior management rejected the recommendation citing custody revenue. The firm's compensation structure rewards relationship managers for custody assets under management without risk adjustment. The board has never reviewed AML metrics for the custody business.

\vspace{0.1cm}
\noindent\textbf{Correct answer approach:} The correct answer will state that custody revenue cannot justify retaining a high-risk customer. The firm must have independent AML decision-making and file SARs regardless of revenue impact. Eliminate answers that suggest custody relationships have different AML standards than other relationships. Remember: revenue is never a valid reason to overlook AML obligations. Senior management override of compliance for revenue is a culture failure.
\end{frame}

% ============================================================
% SLIDE 11 - TRADING LOSSES AS LAYERING
% ============================================================
\begin{frame}{11. Trading Losses as Money Laundering Layering}
\tiny
\noindent\textbf{Introduction to the topic:} Criminals sometimes accept trading losses to clean money. A pattern of consistent small losses across many trades — especially in illiquid securities or complex derivatives — may indicate layering. The criminal deposits illicit funds, executes a series of losing trades, and then withdraws the reduced but now “clean” funds. The losses are the laundering cost. Legitimate traders seek to minimize losses; criminals accept losses as a cost of business.

\vspace{0.1cm}
\noindent\textbf{Type of violation or warning signs:} A customer consistently loses money on trades (e.g., 80\% of trades are losers) but continues trading with the same pattern. The customer cannot explain their trading strategy. The losses are consistent in percentage (e.g., always 2-5\% per trade). The customer trades illiquid securities where prices can be manipulated. The customer trades through multiple accounts across multiple brokers with the same losing pattern. Withdrawals occur after a series of losses, not after gains.

\vspace{0.1cm}
\noindent\textbf{Correct answer approach:} The correct answer will identify that consistent trading losses without a legitimate strategy indicate potential layering. For exam questions, look for answers that require the broker to ask for the customer's trading strategy and file SAR if the explanation is inadequate or the pattern continues. Eliminate answers that assume all trading losses are legitimate business losses. Remember: the pattern matters — consistent losses across many trades is not rational investment behavior.
\end{frame}

% ============================================================
% SLIDE 12 - RAPID IN-AND-OUT TRADING
% ============================================================
\begin{frame}{12. Rapid In-and-Out Trading (Day Trading)}
\tiny
\noindent\textbf{Introduction to the topic:} Rapid in-and-out trading (multiple trades in the same security within a single day) is legitimate for day traders. However, criminals use rapid trading to create transaction volume that obscures the origin of funds. The pattern of rapid trading with no net profit or consistent small losses is suspicious. Brokers must distinguish between legitimate day trading (which has a strategy and produces profits or at least explanations) and layering disguised as day trading.

\vspace{0.1cm}
\noindent\textbf{Type of violation or warning signs:} The customer executes hundreds of trades daily but has no net profit over time. The customer cannot explain their trading strategy when asked. The customer trades the same security repeatedly, buying and selling within minutes. The customer's declared income cannot support the trading volume (e.g., \$50,000 income but \$10 million in daily trading). The customer uses multiple brokerage accounts to fragment trading. The customer withdraws funds after periods of high volume but no net profit.

\vspace{0.1cm}
\noindent\textbf{Correct answer approach:} The correct answer will state that rapid trading without a legitimate strategy or without net profit may be layering and requires enhanced monitoring or SAR filing. For exam questions, look for answers that require the broker to understand the customer's trading strategy and verify that the customer has the financial means to support the trading volume. Eliminate answers that assume all day trading is legitimate. Remember: legitimate day traders have a strategy and can explain it.
\end{frame}

% ============================================================
% SLIDE 13 - WASH TRADING AND MARKING THE CLOSE
% ============================================================
\begin{frame}{13. Wash Trading and Marking the Close}
\tiny
\noindent\textbf{Introduction to the topic:} Wash trading involves buying and selling the same security at the same price, creating artificial volume with no change in beneficial ownership. Marking the close involves buying or selling at the end of the trading day to manipulate the closing price. Both are market manipulation techniques and predicate offenses for money laundering. Criminals use wash trading to create false trading history (making funds appear to come from legitimate trading) and to generate tax losses.

\vspace{0.1cm}
\noindent\textbf{Type of violation or warning signs:} The same customer buys and sells the same security within seconds or minutes at the same price. Two accounts with common control trade the same security back and forth. The customer places large orders in the last few minutes of trading to move the closing price. The customer has no economic reason for the trades (no price movement, no profit). The customer cannot explain the purpose of the trades. The customer requests that trades be executed at specific prices that do not reflect the market.

\vspace{0.1cm}
\noindent\textbf{Correct answer approach:} The correct answer will state that wash trading and marking the close are market manipulation and money laundering predicate offenses requiring SAR filing. For exam questions, look for answers that require the broker to detect these patterns through trade surveillance and report them. Eliminate answers that suggest these are only exchange rule violations, not AML concerns. Remember: the artificial trading activity creates a false appearance of legitimate trading to clean money.
\end{frame}

% ============================================================
% SLIDE 14 - BROKERAGE ACCOUNTS FOR OFFSHORE TRUSTS
% ============================================================
\begin{frame}{14. Brokerage Accounts for Offshore Trusts}
\tiny
\noindent\textbf{Introduction to the topic:} Offshore trusts are legitimate estate planning tools, but they are also used to hide beneficial ownership of brokerage accounts. The trust may have a corporate trustee in a secrecy jurisdiction, a protector with veto power, and discretionary beneficiaries. The broker may not know who the true beneficiaries are. Without identification of all beneficiaries, the broker cannot assess PEP exposure, sanctions risk, or source of funds.

\vspace{0.1cm}
\noindent\textbf{Type of violation or warning signs:} The trust refuses to provide beneficiary names, citing privacy laws. The trust has a flee clause allowing asset movement if threatened. The trust has a letter of wishes that the trustee refuses to share. The trust has been migrated from one jurisdiction to another multiple times. The trust has anonymous beneficiaries or a “class of family members” with no names. The trust engages in unusual trading activity (rapid in-and-out, wash trading) that generates losses.

\vspace{0.1cm}
\noindent\textbf{Correct answer approach:} The correct answer will state that brokers must identify all trust beneficiaries, regardless of discretionary provisions, before opening an account. If the trust refuses to provide beneficiary names, the broker must reject the account. Eliminate answers that suggest trusts are exempt from beneficial ownership identification or that the trustee's representation is sufficient. Remember: FATF requires identification of all beneficiaries; no exceptions for discretionary trusts.
\end{frame}

% ============================================================
% SLIDE 15 - SURVEILLANCE FOR CROSS-PRODUCT MANIPULATION
% ============================================================
\begin{frame}{15. Cross-Product Manipulation (Equities + Derivatives)}
\tiny
\noindent\textbf{Introduction to the topic:} Criminals manipulate prices in one product (e.g., an illiquid stock) to profit in a related product (e.g., options on that stock). This cross-product manipulation is harder to detect because the manipulation and the profit occur in different accounts or different products. The proceeds from the derivative trade are laundered through the brokerage account. Brokers must have surveillance that looks across related products, not just within each product individually.

\vspace{0.1cm}
\noindent\textbf{Type of violation or warning signs:} A customer buys large blocks of an illiquid stock, driving up the price, then sells call options on that stock before the price declines. A customer with no history of derivatives trading suddenly has highly profitable option trades. The underlying stock has unusual volume or price movement with no news. The customer's stock trading loses money but option trading profits. The customer uses different accounts for the stock trades and the derivative trades. The customer cannot explain the relationship between the trades.

\vspace{0.1cm}
\noindent\textbf{Correct answer approach:} The correct answer will state that brokers must have surveillance that detects cross-product manipulation and that such manipulation is a predicate offense for money laundering requiring SAR filing. For exam questions, look for answers that require integration of surveillance across product types and account linkages. Eliminate answers that suggest isolated product surveillance is sufficient. Remember: the manipulation and the profit may be in different accounts controlled by the same person.
\end{frame}

% ============================================================
% SUMMARY SLIDE - BROKER-DEALER AML PRINCIPLES
% ============================================================
\begin{frame}{Summary: Key Broker-Dealer AML Principles}
\tiny
\noindent\textbf{How to answer broker-dealer AML exam questions:}

\vspace{0.1cm}
\noindent• Broker-dealers have full AML obligations under BSA and FINRA Rule 3310. No exemptions for non-custody or small firms.

\vspace{0.1cm}
\noindent• Omnibus accounts require customer identification agreements. If introducing broker refuses to provide underlying customer names, terminate.

\vspace{0.1cm}
\noindent• Insider trading and market manipulation are predicate offenses for money laundering. File SARs for these violations.

\vspace{0.1cm}
\noindent• Cross-border flows without economic purpose indicate layering. Three or more hops without investment activity = SAR.

\vspace{0.1cm}
\noindent• Prime brokers must identify underlying hedge fund investors. Refusal = decline relationship.

\vspace{0.1cm}
\noindent• Sponsored access without pre-trade risk controls is prohibited or requires significant mitigation.

\vspace{0.1cm}
\noindent• IPO and PIPE investors must provide verifiable source of funds. Shell company structures require EDD.

\vspace{0.1cm}
\noindent• Trading losses as a consistent pattern indicate layering. Legitimate traders seek profits; criminals accept losses as cost.

\vspace{0.1cm}
\noindent• Rapid in-and-out trading without a strategy or without net profit may be layering. Ask for strategy explanation.

\vspace{0.1cm}
\noindent• Wash trading and marking the close are manipulation and money laundering predicate offenses.

\vspace{0.1cm}
\noindent• Offshore trust brokerage accounts require identification of all beneficiaries. Refusal = reject.

\vspace{0.1cm}
\noindent• Cross-product manipulation requires integrated surveillance. Profit in derivatives + loss in stocks = red flag.

\vspace{0.1cm}
\noindent• Revenue from custody or trading commissions is never a valid reason to retain a high-risk customer.

\vspace{0.3cm}
\centering
{Remember: Brokers execute trades — they are gatekeepers. When in doubt, file SAR.}
\end{frame}

\end{document}